% \documentclass[journal]{IEEEtran}
\documentclass[journal, onecolumn]{IEEEtran}

\usepackage[T1]{fontenc}
\usepackage{amsmath,amsfonts}
\usepackage{newtxtext}
\usepackage{newtxmath}

\usepackage{algorithmic}

\usepackage{cite}
\usepackage{graphicx}
\usepackage{booktabs}

\usepackage[hidelinks]{hyperref}

\begin{document}

\title{Circuit Diagram for a Wearable Anxiety Monitor}

\author{Andy Babcock \\
    \textit{Steve Sanghi College of Engineering, Northern Arizona University} \\
    Flagstaff, AZ, USA \\
    aab726@nau.edu
}

\maketitle
To realize a wearable anxiety monitor, a NeuroSky MindWave Mobile 2 is being used. This headset was released in 2018, and included first-party software support on Windows, Android, and iOS~\cite{neurosky2018mwm2}. However, it appears that those pieces of software didn't receive significant updates post-release. Though the MindWave Mobile 2 is a relatively cheap electroencephalogram (EEG) headset, there are optimized algorithms that allow for accurate anxiety tracking, even on cheap EEGs~\cite{ChatterjeeDebatri2023Doms}.

For the physical wiring, the MindWave Mobile 2 only supports Bluetooth, with about equal software support on Windows and Android (the two OS's used by myself and Cameron). While direct sync to a PC or laptop using NeuroSky's software would likely work fine, I only have Linux installed on my laptop and desktop, potentially requiring a re-implementation of NeuroSky's software in a third-party program like OpenViBE~\cite{neurosky2018mwm2,openvibe2013}. However, both of us have Android phones, making NeuroSky's eegID mobile app a better option~\cite{neurosky2026eegid}. Setting up this app is relatively easy, requiring a manual connection in Bluetooth settings, then the app handles data recording and processing. To sync that data from the phone to a computer, it is easiest to export as a CSV file post-hoc. However, it would also likely be possible to stream the collected data in real-time or close to real-time using UDP or TCP (UDP is faster with fewer integrity checks on the data transmitted), though it may require a third-party app instead of eegID.

\begin{figure}[htbp]
    \centering
    \includegraphics[width=\columnwidth]{diagram.png}
    \caption{Proposed System Diagram}
    \label{fig:my_image}
\end{figure}

\bibliographystyle{IEEEtran}
\bibliography{references}

\end{document}
